\documentclass{article}
% ======================================================================
% EMBEDDED GE2340 TEMPLATE STYLE (preprint): editable here, no .sty file needed
% Original attribution retained from ge2340_project.sty.
% ======================================================================
% GE2340 course-project style, adapted from the NeurIPS 2024 style:
%
% - uses more LaTeX conventions
% - line numbers at submission time replaced with aligned numbers from
%   lineno package
% - provides [final], [preprint], and [nonatbib] package options
% - automatically loads times package for authors
% - loads natbib automatically; this can be suppressed with the
%   [nonatbib] package option
% - supports named group submissions for the GE2340 course project
% - adds a GE2340 notice in final/submission mode
%
% Original style authors include Roman Garnett and contributors to the
% NeurIPS submission style. Course-specific adaptation by the GE2340 team.
%
% last course revision: September 2026

\makeatletter
\newif\if@gefinal\@gefinalfalse
\newif\if@natbib\@natbibtrue
\newif\if@preprint\@preprinttrue
% determine whether this is an anonymized submission
\newif\if@submission\@submissiontrue
\if@gefinal\@submissionfalse\fi
\if@preprint\@submissionfalse\fi

% fonts
\renewcommand{\rmdefault}{ptm}
\renewcommand{\sfdefault}{phv}

% Course metadata used by the notice string at the bottom of the first page.
\newcommand{\@gecoursecode}{GE2340}
\newcommand{\@geacademicyear}{2026/2027}
\newcommand{\@gesemester}{Semester A}

% acknowledgments
\usepackage{environ}
\newcommand{\acksection}{\section*{Acknowledgments and Disclosure of Funding}}
\NewEnviron{ack}{%
  \acksection
  \BODY
}


% load natbib unless told otherwise
\if@natbib
  \RequirePackage{natbib}
\fi

% set page geometry
\usepackage[verbose=true,letterpaper]{geometry}
\AtBeginDocument{
  \newgeometry{
    textheight=9in,
    textwidth=5.5in,
    top=1in,
    headheight=12pt,
    headsep=25pt,
    footskip=30pt
  }
  \@ifpackageloaded{fullpage}
    {\PackageWarning{ge2340_project}{fullpage package not allowed! Overwriting formatting.}}
    {}
}

\widowpenalty=10000
\clubpenalty=10000
\flushbottom
\sloppy


% font sizes with reduced leading
\renewcommand{\normalsize}{%
  \@setfontsize\normalsize\@xpt\@xipt
  \abovedisplayskip      7\p@ \@plus 2\p@ \@minus 5\p@
  \abovedisplayshortskip \z@ \@plus 3\p@
  \belowdisplayskip      \abovedisplayskip
  \belowdisplayshortskip 4\p@ \@plus 3\p@ \@minus 3\p@
}
\normalsize
\renewcommand{\small}{%
  \@setfontsize\small\@ixpt\@xpt
  \abovedisplayskip      6\p@ \@plus 1.5\p@ \@minus 4\p@
  \abovedisplayshortskip \z@  \@plus 2\p@
  \belowdisplayskip      \abovedisplayskip
  \belowdisplayshortskip 3\p@ \@plus 2\p@   \@minus 2\p@
}
\renewcommand{\footnotesize}{\@setfontsize\footnotesize\@ixpt\@xpt}
\renewcommand{\scriptsize}{\@setfontsize\scriptsize\@viipt\@viiipt}
\renewcommand{\tiny}{\@setfontsize\tiny\@vipt\@viipt}
\renewcommand{\large}{\@setfontsize\large\@xiipt{14}}
\renewcommand{\Large}{\@setfontsize\Large\@xivpt{16}}
\renewcommand{\LARGE}{\@setfontsize\LARGE\@xviipt{20}}
\renewcommand{\huge}{\@setfontsize\huge\@xxpt{23}}
\renewcommand{\Huge}{\@setfontsize\Huge\@xxvpt{28}}

% sections with less space
\providecommand{\section}{}
\renewcommand{\section}{%
  \@startsection{section}{1}{\z@}%
                {-2.0ex \@plus -0.5ex \@minus -0.2ex}%
                { 1.5ex \@plus  0.3ex \@minus  0.2ex}%
                {\large\bf\raggedright}%
}
\providecommand{\subsection}{}
\renewcommand{\subsection}{%
  \@startsection{subsection}{2}{\z@}%
                {-1.8ex \@plus -0.5ex \@minus -0.2ex}%
                { 0.8ex \@plus  0.2ex}%
                {\normalsize\bf\raggedright}%
}
\providecommand{\subsubsection}{}
\renewcommand{\subsubsection}{%
  \@startsection{subsubsection}{3}{\z@}%
                {-1.5ex \@plus -0.5ex \@minus -0.2ex}%
                { 0.5ex \@plus  0.2ex}%
                {\normalsize\bf\raggedright}%
}
\providecommand{\paragraph}{}
\renewcommand{\paragraph}{%
  \@startsection{paragraph}{4}{\z@}%
                {1.5ex \@plus 0.5ex \@minus 0.2ex}%
                {-1em}%
                {\normalsize\bf}%
}
\providecommand{\subparagraph}{}
\renewcommand{\subparagraph}{%
  \@startsection{subparagraph}{5}{\z@}%
                {1.5ex \@plus 0.5ex \@minus 0.2ex}%
                {-1em}%
                {\normalsize\bf}%
}
\providecommand{\subsubsubsection}{}
\renewcommand{\subsubsubsection}{%
  \vskip5pt{\noindent\normalsize\rm\raggedright}%
}

% float placement
\renewcommand{\topfraction      }{0.85}
\renewcommand{\bottomfraction   }{0.4}
\renewcommand{\textfraction     }{0.1}
\renewcommand{\floatpagefraction}{0.7}

\newlength{\@geabovecaptionskip}\setlength{\@geabovecaptionskip}{7\p@}
\newlength{\@gebelowcaptionskip}\setlength{\@gebelowcaptionskip}{\z@}

\setlength{\abovecaptionskip}{\@geabovecaptionskip}
\setlength{\belowcaptionskip}{\@gebelowcaptionskip}

% swap above/belowcaptionskip lengths for tables
\renewenvironment{table}
  {\setlength{\abovecaptionskip}{\@gebelowcaptionskip}%
   \setlength{\belowcaptionskip}{\@geabovecaptionskip}%
   \@float{table}}
  {\end@float}

% footnote formatting
\setlength{\footnotesep }{6.65\p@}
\setlength{\skip\footins}{9\p@ \@plus 4\p@ \@minus 2\p@}
\renewcommand{\footnoterule}{\kern-3\p@ \hrule width 12pc \kern 2.6\p@}
\setcounter{footnote}{0}

% paragraph formatting
\setlength{\parindent}{\z@}
\setlength{\parskip  }{5.5\p@}

% list formatting
\setlength{\topsep       }{4\p@ \@plus 1\p@   \@minus 2\p@}
\setlength{\partopsep    }{1\p@ \@plus 0.5\p@ \@minus 0.5\p@}
\setlength{\itemsep      }{2\p@ \@plus 1\p@   \@minus 0.5\p@}
\setlength{\parsep       }{2\p@ \@plus 1\p@   \@minus 0.5\p@}
\setlength{\leftmargin   }{3pc}
\setlength{\leftmargini  }{\leftmargin}
\setlength{\leftmarginii }{2em}
\setlength{\leftmarginiii}{1.5em}
\setlength{\leftmarginiv }{1.0em}
\setlength{\leftmarginv  }{0.5em}
\def\@listi  {\leftmargin\leftmargini}
\def\@listii {\leftmargin\leftmarginii
              \labelwidth\leftmarginii
              \advance\labelwidth-\labelsep
              \topsep  2\p@ \@plus 1\p@    \@minus 0.5\p@
              \parsep  1\p@ \@plus 0.5\p@ \@minus 0.5\p@
              \itemsep \parsep}
\def\@listiii{\leftmargin\leftmarginiii
              \labelwidth\leftmarginiii
              \advance\labelwidth-\labelsep
              \topsep    1\p@ \@plus 0.5\p@ \@minus 0.5\p@
              \parsep    \z@
              \partopsep 0.5\p@ \@plus 0\p@ \@minus 0.5\p@
              \itemsep \topsep}
\def\@listiv {\leftmargin\leftmarginiv
              \labelwidth\leftmarginiv
              \advance\labelwidth-\labelsep}
\def\@listv  {\leftmargin\leftmarginv
              \labelwidth\leftmarginv
              \advance\labelwidth-\labelsep}
\def\@listvi {\leftmargin\leftmarginvi
              \labelwidth\leftmarginvi
              \advance\labelwidth-\labelsep}

% create title
\providecommand{\maketitle}{}
\renewcommand{\maketitle}{%
  \par
  \begingroup
    \renewcommand{\thefootnote}{\fnsymbol{footnote}}
    % for perfect author name centering
    \renewcommand{\@makefnmark}{\hbox to \z@{$^{\@thefnmark}$\hss}}
    % The footnote-mark was overlapping the footnote-text,
    % added the following to fix this problem               (MK)
    \long\def\@makefntext##1{%
      \parindent 1em\noindent
      \hbox to 1.8em{\hss $\m@th ^{\@thefnmark}$}##1
    }
    \thispagestyle{empty}
    \@maketitle
    \@thanks
    \@notice
  \endgroup
  \let\maketitle\relax
  \let\thanks\relax
}

% rules for title box at top of first page
\newcommand{\@toptitlebar}{
  \hrule height 4\p@
  \vskip 0.25in
  \vskip -\parskip%
}
\newcommand{\@bottomtitlebar}{
  \vskip 0.29in
  \vskip -\parskip
  \hrule height 1\p@
  \vskip 0.09in%
}

% create title (includes both anonymized and non-anonymized versions)
\providecommand{\@maketitle}{}
\renewcommand{\@maketitle}{%
  \vbox{%
    \hsize\textwidth
    \linewidth\hsize
    \vskip 0.1in
    \@toptitlebar
    \centering
    {\LARGE\bf \@title\par}
    \@bottomtitlebar
    \if@submission
      \begin{tabular}[t]{c}\bf\rule{\z@}{24\p@}
        Anonymous Author(s) \\
        Affiliation \\
        Address \\
        \texttt{email} \\
      \end{tabular}%
    \else
      \def\And{%
        \end{tabular}\hfil\linebreak[0]\hfil%
        \begin{tabular}[t]{c}\bf\rule{\z@}{24\p@}\ignorespaces%
      }
      \def\AND{%
        \end{tabular}\hfil\linebreak[4]\hfil%
        \begin{tabular}[t]{c}\bf\rule{\z@}{24\p@}\ignorespaces%
      }
      \begin{tabular}[t]{c}\normalfont\rule{\z@}{24\p@}\@author\end{tabular}%
    \fi
    \vskip 0.3in \@minus 0.1in
  }
}

% add the course notice to the bottom of the first page
\newcommand{\ftype@noticebox}{8}
\newcommand{\@notice}{%
  % give a bit of extra room back to authors on first page
  \enlargethispage{2\baselineskip}%
  \@float{noticebox}[b]%
    \footnotesize\@noticestring%
  \end@float%
}

% abstract styling
\renewenvironment{abstract}%
{%
  \vskip 0.075in%
  \centerline%
  {\large\bf Abstract}%
  \vspace{0.5ex}%
  \begin{quote}%
}
{
  \par%
  \end{quote}%
  \vskip 1ex%
}

% For the paper checklist
\newcommand{\answerYes}[1][]{\textcolor{blue}{[Yes] #1}}
\newcommand{\answerNo}[1][]{\textcolor{orange}{[No] #1}}
\newcommand{\answerNA}[1][]{\textcolor{gray}{[NA] #1}}
\newcommand{\answerTODO}[1][]{\textcolor{red}{\bf [TODO]}}
\newcommand{\justificationTODO}[1][]{\textcolor{red}{\bf [TODO]}}

% handle tweaks for camera-ready copy vs. submission copy
\if@preprint
  \newcommand{\@noticestring}{%
    % Preprint. Under review.%
  }
\else
  \if@gefinal
    \newcommand{\@noticestring}{%
      \@gecoursecode\/ Course Project, \@geacademicyear\/ \@gesemester.%
    }
  \else
    \newcommand{\@noticestring}{%
      Submitted to \@gecoursecode\/ Course Project,
      \@geacademicyear\/ \@gesemester. Internal use only.%
    }

    % hide the acknowledgements
    \NewEnviron{hide}{}
    \let\ack\hide
    \let\endack\endhide

    % line numbers for submission
    \RequirePackage{lineno}
    \linenumbers

    % fix incompatibilities between lineno and amsmath, if required, by
    % transparently wrapping linenomath environments around amsmath
    % environments
    \AtBeginDocument{%
      \@ifpackageloaded{amsmath}{%
        \newcommand*\patchAmsMathEnvironmentForLineno[1]{%
          \expandafter\let\csname old#1\expandafter\endcsname\csname #1\endcsname
          \expandafter\let\csname oldend#1\expandafter\endcsname\csname end#1\endcsname
          \renewenvironment{#1}%
                          {\linenomath\csname old#1\endcsname}%
                          {\csname oldend#1\endcsname\endlinenomath}%
        }%
        \newcommand*\patchBothAmsMathEnvironmentsForLineno[1]{%
          \patchAmsMathEnvironmentForLineno{#1}%
          \patchAmsMathEnvironmentForLineno{#1*}%
        }%
        \patchBothAmsMathEnvironmentsForLineno{equation}%
        \patchBothAmsMathEnvironmentsForLineno{align}%
        \patchBothAmsMathEnvironmentsForLineno{flalign}%
        \patchBothAmsMathEnvironmentsForLineno{alignat}%
        \patchBothAmsMathEnvironmentsForLineno{gather}%
        \patchBothAmsMathEnvironmentsForLineno{multline}%
      }
      {}
    }
  \fi
\fi



\makeatother
% END EMBEDDED TEMPLATE STYLE

\usepackage[utf8]{inputenc}
\usepackage[T1]{fontenc}
\usepackage{url}
\usepackage{booktabs}
\usepackage{longtable}
\usepackage{array}
\newcolumntype{L}[1]{>{\raggedright\arraybackslash}p{#1}}
\usepackage{amsmath,amssymb}
\usepackage{graphicx}
\usepackage{tikz}
\usetikzlibrary{arrows.meta}
\usepackage{microtype}
\usepackage{xcolor}
\usepackage{float}
\usepackage{hyperref}
\setcitestyle{numbers,square}
\hypersetup{colorlinks=true,linkcolor=blue,citecolor=blue,urlcolor=blue,
  pdftitle={Phase 2: Course Registration Helper for Balanced Workloads},pdfauthor={}}
\urlstyle{same}
\title{Phase 2: AI Project Report\\
\large Course Registration Helper for Balanced Workloads}
\author{
  \underline{\hspace{1in}} (Student ID: \underline{\hspace{0.65in}})\\
  \texttt{\underline{\hspace{2.2in}}}
  \And
  \underline{\hspace{1in}} (Student ID: \underline{\hspace{0.65in}})\\
  \texttt{\underline{\hspace{2.2in}}}
  \AND
  \underline{\hspace{1in}} (Student ID: \underline{\hspace{0.65in}})\\
  \texttt{\underline{\hspace{2.2in}}}
  \And
  \underline{\hspace{1in}} (Student ID: \underline{\hspace{0.65in}})\\
  \texttt{\underline{\hspace{2.2in}}}
}


\begin{document}
\maketitle
\section{Introduction \& Problem Recap}
Study-path planning is an important progress in university life.However, for some students in city university unfamiliar with GE or elective courses ,such as its background requirements, workload. While seats are limited, students often have to make a rush decision for which course to register. And there is a more often appears that: originally there is no seats at the before the semester start. After the first week, some other students in certain dropped the course, which free up the seats. However, the students who failed to register any GE courses ,rush to register the courses with new free space. But they may get into a course which difficult to them at week 2, where the add-drop period has been ended, make the decision become no no-regret.
Our helper uses an LLM and course scoring to suggest suitable, workload-aware schedules from students' interests, year and required courses. The prototype covers selected courses in one semester; it supports course selection rather than automatic registration or guaranteed seats.

Our design combines a fine-tuned language model with mathematical optimisation. The model interprets course information and proposes structured recommendations; source checks and an optimiser turn them into a constrained plan. This lets us personalise recommendations while respecting programme eligibility and academic dependencies.



\section{Data Collection \& Methodology}
\textbf{Programme-specific evidence.} We assembled the handbook, CM 2025-entry path/course list, CS major/suggested-plan documents, GE information and saved Semester A/B timetables\citep{cityu-programmes}. We extract identities, credits, prerequisites, exclusions, restrictions and meetings, retaining source cohorts and years. Courses must appear in the student's designated programme document or verified GE collection. and the elective courses is labeled just for CS major or CM major,while GE courses are in common pool for selection 


\textbf{PDFs data used for training} We collected 93 distinct timetable PDFs: 86 parsed successfully and seven were excluded. A split-level audit identifies 82 timetable PDFs and three programme PDFs represented in the 485 training examples: \textbf{85 PDF sources in this programme-training split}. The three programme sources are the CM course list, CS major requirements and CS suggested plan. We teach extracted text/facts as JSON examples, not raw PDF files. Earlier handbook, GE and study-path stages are outside this scoped count.

\subsection{Model adaptation and prompts.} We adapted Qwen3-8B using rank-8 QLoRA, a frozen four-bit base and BF16 computation\citep{qwen,qlora}. The programme dataset contains 485/24/30 training/validation/test examples; the CM background update uses 244/32 plus 20 held-out checks. Relation distillation uses 149/20 examples, followed by an 82-example maximum-versus-sum correction. These evaluate familiar evidence and synthetic policies, not real workload or unseen-cohort generalisation.

Training uses system instructions, user-supplied facts/states and assistant JSON targets, with assistant-only cross-entropy loss. The current prompt permits only \texttt{reward\_parameters} and \texttt{course\_relations}. One local inference exported all eight settings and 26 links without fallback. This checks structured policy export, not independent reward quality. The latest optimiser revision did not retrain neural weights.
\begingroup\setlength{\parskip}{3pt}\setlength{\intextsep}{6pt}
\subsection{Plan for training GPT-6 sol teacher and human feedback training} Figure~\ref{fig:workflow} shows the implemented division of work. In the beginning stages, we use GPT-6 sol as a teacher to train the small Qwen model. The GPT will use it deductive ability and make Qwen learn the reward based on the course relation in certain study path. After the GPT- training round, we start generate fictional examples to see the results, and give human feedback to the model for the retraining, or rules define if the results is impossible to be achieve in reality.\\
\begin{figure}[H]
\centering
% EDITABLE WORKFLOW: edit the text inside each node below.
% All node positions are in points. Human review follows optimizer output.
% No external PDF or image file is needed for this diagram.
\resizebox{\linewidth}{!}{%
\begin{tikzpicture}[
  x=1pt,y=1pt,
  flowbox/.style={draw=blue!35!black,fill=blue!4,rounded corners=3pt,
    line width=0.6pt,align=center,text width=82pt,minimum height=34pt,
    inner sep=3pt,font=\sffamily\fontsize{8}{9}\selectfont},
  flowarrow/.style={-{Latex[length=4pt,width=3pt]},draw=blue!35!black,
    line width=0.6pt}
]
  % Main workflow: GPT teaches the reward policy; the tool selects the path.
  \node[flowbox] (teacher) at (47,117)
    {\textbf{GPT teacher}\\Reward policy\\Course relations};
  \node[flowbox] (qwen) at (149,117)
    {\textbf{Qwen reward model}\\8 settings, 26 links\\Structured JSON};
  \node[flowbox] (checker) at (251,117)
    {\textbf{Parameter checker}\\Schema and bounds\\Rules stay fixed};
  \node[flowbox] (optimizer) at (353,117)
    {\textbf{External optimizer}\\Exact DP search\\Path + certificate};
  \draw[flowarrow] (teacher.east) -- (qwen.west);
  \draw[flowarrow] (qwen.east) -- (checker.west);
  \draw[flowarrow] (checker.east) -- (optimizer.west);
  % Optimizer output is reviewed by a human before feedback is applied.
  \node[flowbox] (human) at (353,67)
    {\textbf{Human review}\\Inspect output\\Give advice};
  \node[flowbox] (feedback) at (149,17)
    {\textbf{Human feedback}\\Update parameters\\and planning rules};
  \node[flowbox,text width=184pt] (facts) at (302,17)
    {\textbf{Student profile + official course facts}\\
     Credits, prerequisites and A/B availability\\
     Saved meetings / pending timetable flags};
  \draw[flowarrow] (feedback.north) -- (qwen.south);
  \draw[flowarrow] (optimizer.south) -- (human.north);
  \draw[flowarrow] (human.west) -- (199,67) -- (199,17) -- (feedback.east);
  \draw[flowarrow] (facts.east) -- (411,17) -- (411,117) -- (optimizer.east);
\end{tikzpicture}
}
\caption{The optimizer sends its output to human review. Human feedback advises updates to reward parameters and planning rules, which return through Qwen and validation before optimisation is repeated. Official course facts remain a separate input.}
\label{fig:workflow}
\end{figure}

\subsection{Why LLM teacher and why Qwen as input ?} We use GPT-6 sol as a based teacher in the beginning stage to reduce the labor effort for simple correlation between course and course. It can help the model to decide what is best suitable for students, given they have some knowledge on some courses from transcript.\\
And the reasons that we train a Qwen and release the parameter from qwen is simple: we want the input become more user-friendly and not limited, which the auto-encoder for a small LLM is best for our choices for short time black box method. However, we believe the Qwen part can be replaced by more simipler encoder, and we will discuss it in the discussion section.


\section{Implementation Details}
\subsection{Hardware and GPU memory usage}
We trained and ran Qwen locally on a single NVIDIA GeForce RTX 5090. CUDA reports 31.82\,GiB of total device memory. The relevant training runs use rank-8 QLoRA with a frozen four-bit NF4 base and BF16 computation. Training processes one example at a time and accumulates gradients over four examples. The programme and relation runs allow sequences of up to 4,096 tokens; the degree-policy run allows up to 6,144 tokens.

Recorded peak PyTorch-allocated GPU memory was 8.27\,GiB for programme training, 7.37\,GiB for the finance-background update, 8.35\,GiB for relation correction, and 15.84\,GiB for the longer-context degree-policy update. These are measurements from separate runs, not memory used simultaneously.

A representative inference measurement on 26 November 2026 used the four-bit relation model, batch size one, a 598-token prompt and a 33-token response for a course-preparation reward query. Model-loaded allocation was 5.74\,GiB; peak inference allocation was 5.99\,GiB and peak reserved memory was 6.10\,GiB. These figures cover this process's PyTorch GPU memory, including model and generation, and exclude other applications and driver memory. Longer prompts, responses or larger batches can change consumption. The subsequent Python path search runs on the CPU using the validated reward parameters.


\subsection{User input and personalisation}
The student supplies their programme (CM or CS), entry cohort, current study year and semester, and completed-course history. The prototype represents these inputs as a structured student profile containing course codes, credits, grades and completed terms. To further simplify, the model accept a transcript as a direct input, which completely cover those inputs mentioned in above. 

To further personalize the study path, user can also specify a desired stream and interests to guide the recommended path. A desired stream increases the reward for progress towards that stream. Also, there are 4 adjustable non-negative weights $\lambda_b$, $\lambda_d$ and $\lambda_w$ to express how strongly the planner should penalise background mismatch, difficulty and workload.(Details will be at section 3.4 and 3.5)  

\subsection{Basic rules for model:graduation and evidence constraints.} The terminal audit requires at least 122 unique credits (therefore also satisfying 121), four distributional GEs covering all three areas, at least 20 elective mathematics credits outside core/capstone study, mandatory/core completion, the project/free-elective requirements and the chosen stream. At most one distributional GE is selected per semester; compulsory language/development GE is separate. Loads are flexible, with zero to six positive-credit courses and at most 18 credits per term.
\subsection{Basic penalty and reward function} started with minimised weighted penalty function $P(x)$ for binary course choices $x_i$. 
\[\min_{\mathbf x\in\mathcal X}P(\mathbf x),\quad P(\mathbf x)=\sum_i x_i(\lambda_b b_i+\lambda_d d_i+\lambda_w w_i)\]
Here $b,d,w$ denote background mismatch, difficulty and workload, with non-negative student-adjustable $\lambda$ weights. Higher cost lowers reward; degree constraints $\mathcal X$ remain binding. \\
Now consider a reward $R_i$ at choice $x_i$ , we can denote the objective reward function  
\[\max_{\mathbf x\in\mathcal X}\left[\sum_i x_iR_i\right]\]
combine them we have
\[\max_{\mathbf x\in\mathcal X}\sum_i x_iR_i-P(\mathbf x)\]
and we can restate the function as a State-dependent reward function
\subsection{(important) State-dependent reward function.} Expanding $R_i$ into preparation, graduation and stream progress gives the implemented Markov-state utility. For selected courses $X$, we maximise
\begin{align}
U(X,s)={}&\sum_{c\in X}\bigl[g_c(s)+p_c(s)+h_c(s)+q_c(s)+z_c(s)\nonumber\\
&-\lambda_b(b_c(s)+e_c(s))-\lambda_d d_c(s)-\lambda_w w_c(s)\bigr]
+R_{\mathrm{stream}}(X,s)-O(X,s).\label{eq:utility}
\end{align}

Base value $g$, preparation $p$, core priorities $h$, prerequisite bridges $q$ and quota/timing bonuses $z$ reward progress. The student-adjustable weights $\lambda_b,\lambda_d,\lambda_w$ penalise background mismatch, difficulty and workload. These scores are configured estimates, not measured study hours. $R_{\mathrm{stream}}(X,s)$ are the stream reward, where claiming one selected stream earns 6 once, multiplied by the student's stream weight.

\textbf{over workload penalty function $O(X,s)$}
We consider if we don't set a upperbound of the credit taken in a semster, the model will only output to go for all course. And we don't want the model to advice many course in same semster, we introduce  over workload penalty function $O(X,s)$ ,where
\begin{equation}
O(X,s)=\begin{cases}
0.25[\max(0,C(X)-15)]^2,&\text{Years 1--3},\\
0.25[\max(0,C(X)-12)]^2,&\text{Year 4}.
\end{cases}\label{eq:overload}
\end{equation}
The 18-credit maximum is a separate hard constraint. A four-year recommendation must satisfy the terminal degree audit; a high reward cannot compensate for missing required credits.


\subsection{Solving optimisation problem.} We maximise the reward objective using exact dynamic programming implemented in Python. Starting from the student's completed courses and fixed allocations, the program explores the permitted semester placements of the remaining courses. For each candidate path, it updates the student's state in semester order and calculates the total discounted reward,and identical partial allocations are merged to avoid repeated calculations. The program compares every feasible completion in the supplied choices and returns the path with the largest objective value.

\section{Results \& Evaluation}
\textbf{Evaluation protocol.} We use fictional transcripts and controlled interventions. Both methods receive the same three histories, fixed context. GPT uses no solver/code/browsing tools and sees no optimiser answers or scores; its final response is scored without repair. We corrected an omitted state-dependent cost before the final comparison and excluded pilot outputs. The exact comparator endpoint is unverified, so we label it GPT.


\textbf{Three matched student inputs.} Each fictional student has completed exactly 15 credits in Year 1A and 15 in Year 1B. The common Year 1A passes are MA1300, CS1302, PHY1201 and GE1401; Year 1B passes are MA1301, CHEM1200, GE1501 and GE2401. Each term also includes the profile-specific three-credit GE in Table~\ref{tab:profiles}. The fixed future course context is supplied identically to both methods within each case; it differs across cases, so this is not a causal comparison of GE histories.

\begin{table}[H]
\centering\small
\caption{Completed Year-1 GE profiles and desired stream. All listed courses are assumed passes.}
\begin{tabular}{L{0.45in}L{1.45in}L{1.45in}L{1.45in}}
\toprule
Profile & Year 1A GE & Year 1B GE & Desired stream\\\midrule
F & GE2260, Area 2 & GE1105, Area 1 & Financial Mathematics\\
N & GE1305, Area 3 & GE1113, Area 1 & Financial Mathematics\\
M & GE2134, Area 1 & GE1305, Area 3 & Enriched Mathematics\\
\bottomrule
\end{tabular}\label{tab:profiles}
\end{table}

\textbf{Actual outputs: reward model plus DP versus direct GPT.} Our method uses Qwen's validated parameter/link packet and the external DP. The baseline GPT selects placements directly from the same packet and course facts, without running the optimisation method or seeing its results. The common critic checks both outputs; it performs no repair or fallback for GPT. Table~\ref{tab:outputs} shows every elective decision, omitting clock times. Terms not listed contain the unchanged fixed course context.

{\small\setlength{\tabcolsep}{4pt}\renewcommand{\arraystretch}{1.1}
\begin{table}[H]\centering\small
\caption{Paired outputs for all six missing electives. Both methods produce the same placements in every case.}
\begin{tabular}{L{0.4in}L{2.35in}L{2.35in}}
\toprule
Profile & Qwen reward model + external DP & Direct GPT, without optimiser\\\midrule
F & Y3B: MA4529, MA4537\par Y4A: MA4528\par Y4B: MA4546, MA4550, MA4551 & Y3B: MA4529, MA4537\par Y4A: MA4528\par Y4B: MA4546, MA4550, MA4551\\\addlinespace[4pt]
N & Y3A: MA4528\par Y4B: MA4529, MA4537, MA4546, MA4550, MA4551 & Y3A: MA4528\par Y4B: MA4529, MA4537, MA4546, MA4550, MA4551\\\addlinespace[4pt]
M & Y3B: MA4537, MA4546\par Y4A: MA4528\par Y4B: MA4529, MA4550, MA4551 & Y3B: MA4537, MA4546\par Y4A: MA4528\par Y4B: MA4529, MA4550, MA4551\\\addlinespace[4pt]
\bottomrule\end{tabular}\label{tab:outputs}\end{table}}


Both methods pass the encoded curriculum audit in all three cases and reach the same certified completion optimum. 


\textbf{A complete output, including fixed courses.} Table~\ref{tab:fullpath} expands profile N into its full projected path. GPT and the optimiser return the same six elective placements, so the resulting full paths are identical. The Year-1 history is unchanged, and all future grades are assumed passes for the simulation.

\begin{table}[H]\centering\small
\caption{Profile N: common full output from both methods. A star marks pending class details; zero-credit catch-up requirements are shown explicitly.}
\begin{tabular}{L{0.38in}L{3.75in}rr}\toprule
Term & Selected courses & CUs & Total\\\midrule
Y1A & MA1300, CS1302, PHY1201, GE1401, GE1305 & 15 & 15\\
Y1B & MA1301, CHEM1200, GE1501, GE2401, GE1113 & 15 & 30\\
Y2A & CS2360, CSCI1001$^{*}$ (0), GE1601$^{*}$, GE2260, MA2503, MA2508, MA2509 & 18 & 48\\
Y2B & CS2468, CSCI1002$^{*}$ (0), EF3320, MA2507, MA3511, MA3515, MA3517 & 16 & 64\\
Y3A & GE1105, MA3512, MA3525, MA4528$^{*}$, MA4545 & 15 & 79\\
Y3B & MA2510, MA3514, MA3523, MA3526, MA3530 & 13 & 92\\
Y4A & MA3518, MA3521, MA4530 (A)$^{*}$, MA4542 & 12 & 104\\
Y4B & MA4529$^{*}$, MA4530 (B)$^{*}$, MA4537$^{*}$, MA4546$^{*}$, MA4550$^{*}$, MA4551$^{*}$ & 18 & 122\\

\bottomrule\end{tabular}\label{tab:fullpath}\end{table}

Twenty-five future credits still lack saved class details. They remain pending, with null meetings/CRNs; project parts denote one six-credit year-long course allocated three credits per term. Missing zero-credit soft-skills requirements are planned catch-up tasks requiring confirmation, not invented historical passes. A curriculum-complete result therefore remains distinct from a confirmed registration timetable.

All three repaired paths contain four distributional GE passes across Areas 1--3, 30 elective mathematics credits, 45 core credits and six project credits, with the selected stream complete under the encoded rules. Profile N's advanced B-term electives follow MA2510 in Year 3B and MA3518/MA3521 in Year 4A. Its 18-credit final term therefore incurs the configured nine-point Year-4 overload penalty. A 12- or 13-credit term remains valid because the whole path now reaches 122; we do not force five courses in every term.

\textbf{Parameter sensitivity changes a tool decision.} Holding profile F and its course links fixed, changing only the Year-4 overload coefficient changes the optimal placement (Table~\ref{tab:sensitivity}). This verifies that the optimiser responds to supplied parameters rather than reproducing a fixed language-model answer. The zero-weight run is a counterfactual test; the configured policy remains 0.25.

\begin{table}[H]\centering\small
\caption{Controlled overload intervention, with 122 credits retained.}
\begin{tabular}{lrrr}\toprule
Year-4 coefficient & Year 3B CUs & Year 4B CUs & Total\\\midrule
0.25, configured & 18 & 15 & 122\\
0, counterfactual & 15 & 18 & 122\\\bottomrule
\end{tabular}\label{tab:sensitivity}\end{table}

\textbf{Dependency experiment: how earlier study changes rewards.} Our earlier 12 paired interventions hold year, base profit, score vector and penalty weights fixed while excluding credit-progress and overload effects. All match their predeclared policy effects. Table~\ref{tab:dependency} retains representative relations; they are computational tests of configured mechanisms.

\begin{table}[H]\centering\small
\caption{Representative dependency effects, with other tested components fixed.}
\begin{tabular}{L{2.5in}L{0.75in}L{1.7in}}\toprule
Changed history or preference & Target & Reward effect\\\midrule
Pass MA2508 & MA3515 & Preparation $+0.20$\\
Pass MA2509 & GE2134 & Preparation $+0.15$\\
Covered Area 1 replaces Area 2 & GE1113 & Base reward $-2.70$\\
Pass GE2260; finance selected & MA3521 & Stream reward $+4.00$\\
Pass MA2503; portfolio mathematics & EF3320 & Preparation $+0.20$\\\bottomrule
\end{tabular}\label{tab:dependency}\end{table}

\section{Discussion}
\subsection{Effective of transfer learning} In the results part, we show the powerful of transfer learning by using a high-parameter LLM as a teacher to a LLM which with small-parameter. With the parameter teaching method, we can achieve an basic similar output after fine-tuning, as showing in the result part.
\subsection{Benefits for switching into a local small model,instead of directly using GPT }
Now it will raise a question, why don't we directly use GPT to suggest a plan if we they results similar?\\
\textbf{privacy} Recall the input may take the transcript of the students, and it will results the best performance since it gives details of the background of the students .However, there may be personal information on the transcript.Directly uploading it to GPT, these information will be potentially exposed through internet or exposed through out the GPT model. Using a finetuned model to run in at local can avoid this issues, since everything is computed locally.\\
\textbf{Benefits of update and as a module} Since the course timetable,information will be updated on each years/ semester, it will be difficult if every time user have to feed a hugh numbers of course data to GPT.  However, by finetuning a small model, the process can be reduced to one man training and sharing to others, which is more update-friendly.\\
Furthermore, a small-size model can be run with simple hardware. Some user,or the official of city university can add ane updated finetuned model to an online chat bot, which further reduce the minimum requirement for user to attempt the method.\\
\subsection{Is LLM necessary during inference?}
Now, consider all rewards, workload parameter is released in certain states, and clearly correlation between course and course is in a finite set, theoretically we can put all the parameters and effect from correlation into a big database. Furthermore, we can reduce the possible optimal path into a smaller size of subset, which reduce the computation.\\
For example, consider a case: a student who finished GE2340 will receive extra reward for future AI course, such as machine learning. Then during the path selection, the shortcut GE2340 $\to \dots \to$ machine learning will be mostly likely the optimal option, by greedy approach. From the process of clear greedy algorithm, we can significantly reduce the amounts of possible paths that can be selected as the optimal solution. And if we fixed the input ,instead of accepting  all type of nature language input with a auto-encoder, the whole process can be reduced into a simple script. \\

\subsection{limitation}
\textbf{further training data} In our build of the model, we only use GPT's advice and personal feedback, but not all courses in the training set are taken by our feedback team. Also, there are no data about some external factors, such as professor's requirement, which may affect the difficulty of the course. At last, the workload parameter is still considered as unreliable, further data collection, such as students' feedback on the course, should be collected for further training.\\
\textbf{Limiting factors for reward function} For the reward function, only few factors ,such as correlation or workload is selected to decide the reward. However there maybe some other factors in real life will affect the decision making, such as the venue of the lectures, or how clumsy the schedule is arranged in a single day (i.e, our helper may generate a schedule with lecture from 9 am to 6 pm in one single day, which the students may not want to take so.), the reward function should be updated if there is any meaning factors that can be added in 
\section{Conclusion}
Our helper uses an LLM for rewards and relations, and external dynamic programming for path selection. During the transfer training and human feedback, the finetuned model showed it high potential on giving personalize suggestion on study path and schedule planning. However, there are still some limitations, such as the lack of training data and factors is limited. It is suggested further updated and improvement can be occur in the further to make the helper become more complete.

\begin{thebibliography}{9}
\footnotesize\setlength{\itemsep}{0pt}\setlength{\parsep}{0pt}
\bibitem[Yang et~al.(2025)]{qwen}
A. Yang et al., ``Qwen3 Technical Report,'' arXiv:2505.09388, 2025. \url{https://arxiv.org/abs/2505.09388}.
\bibitem[Dettmers et~al.(2023)]{qlora}
T. Dettmers et al., ``QLoRA: Efficient Finetuning of Quantized LLMs,'' 2023. \url{https://arxiv.org/abs/2305.14314}.
\bibitem[CityUHK(2022--2026)]{cityu-programmes}
CityUHK, supplied Student Handbook 2026--27, CM 2025-entry normative path and designated course list, CS major/suggested-plan documents, GE and saved Semester A/B timetables. Current advanced-course rules: \url{https://www.cityu.edu.hk/catalogue/ug/current/course/MA4537.htm} and corresponding course-code pages.
\end{thebibliography}
\end{document}
